\section{Inverse-density interpolation for marked spin traces}
\label{ten:inverse}
This argument proves a grid moment estimate with the full marked-site entropy cost. Its operator comparison must be performed before the column operator is split into local sectors. The subsequent interpolation treats products of noncommuting density matrices in their original order.
\smallskip
All notation introduced below is local to this section. Traces are unnormalized, logarithms are natural, and a sweep on $2^d$ positions takes $d$ physical shuffles. The conventions are those of Section~\ref{sec:prelim}.

\subsection{The sweep moments}
Use the sweep $B_n$ of Section~\ref{sec:prelim}, with \(n=2^d\ge2\).
For a one-position grid factor set \(B_1=I\), the empty product of
coordinate layers. There is a grid decomposition for \(n=ab\) with
\(a,b\) powers of two. Take \(a\) rows of length \(b\), so the first axes
of the sweep are inside rows and the remaining ones inside columns. Then
\(B_n=C R\), where \(R\) is a tensor product in the row permutation
subgroup of \(a\) copies of \(B_b\), and \(C\) is analogously made from
\(b\) copies of \(B_a\) in the column subgroup.

We use unitary representations. For a diagram (partition) \(\lambda\) of \(m\), write \(D_\lambda=f^\lambda\) for the dimension of the corresponding irreducible of \(S_m\), \(H_\lambda=\log D_\lambda\) (logs natural). Use dimension 1 for the empty diagram; \(f^{\lambda/\omega}\) counts standard tableaux of the indicated skew shape. The moments used in the proof are
\[
 J_m(\lambda,p)=D_\lambda\,{\rm Tr}_\lambda |B_m|^p,\qquad |A|=(A^*A)^{1/2}.
\]
Here and below a trace with diagram subscript is in one copy of the irreducible. We prove bounds for these moments with bounded \(p\). We use standard representation theory of symmetric groups: in particular the branching rule, Schur-Weyl duality, the Littlewood-Richardson (LR) rule, and the hook length formula. Matrix inequalities used below include the Araki-Lieb-Thirring trace inequality (for positive matrices \(L,M\) and \(v\ge1\), \({\rm Tr}(L^{1/2}M L^{1/2})^v\le {\rm Tr}(L^v M^v)\), trace of powers on the left) and Schatten H\"older.

Roughly, at sparse levels we use the small probability of path encounters for all the cards on which there must be nontrivial dependence. At the other levels we propagate trace-moment bounds from the rows and columns. In that argument a spin representation (with two parities so as to cover both long rows and long columns of a diagram), augmented by distinct marks, permits testing for global permutation types with losses that combine well in entropy scale.

\subsection{The optimized hook budget}

The grid estimate will be used with a hook that grows with the deck.
Fix
\[
 \epsilon=\frac1{1000},\qquad c_0=\frac1{100},\qquad
 \theta=\frac12,\qquad Q_m=\lfloor m^{1/16}\rfloor.
\]
For a partition $\xi\vdash m$, define
\[
 G_m(\xi)=\min_{\substack{\alpha\subseteq\xi\\
                         \alpha_{Q_m+1}\le Q_m}}
 \left\{\theta H_\alpha+c_0(m-|\alpha|)\log m
             -(m-|\alpha|)\log\frac{m}{m-|\alpha|}\right\},
\]
where the last term is zero when $|\alpha|=m$.
Call $\xi$ large-level when $m-\xi_1>m^{1-\epsilon}$.
The minimum retains the full, unclipped entropy cost of the deleted
positions; at the other levels we will use a separate sparse estimate.

\begin{proposition}[Bounded exponents for the inverse-density budget]
\label{s:inverse-bounded-exponent}
There are real exponents $p_d\ge2$, with $\sup_{d\ge1}p_d<\infty$,
such that for every $d\ge1$ and every $\lambda\vdash n=2^d$,
\[
 \log J_n(\lambda,p_d)\le
 \begin{cases}
 G_n(\lambda),&n-\lambda_1>n^{1-\epsilon},\\
 0,&n-\lambda_1\le n^{1-\epsilon}.
 \end{cases}
\]
For all sufficiently large $n$, the second bound strengthens to
$\log J_n(\lambda,p_d)\le-10k\log n$ when
$1\le k=n-\lambda_1\le n^{1-\epsilon}$.
\end{proposition}

The sparse path estimate proves the second branch. For the first branch,
we prove a marked-grid estimate at a prescribed parent hook, extend the
child budgets to that inherited hook, and absorb the resulting error by
a bounded increase in the Schatten exponent. The grid estimate is also
the fixed-parameter consequence of
Corollary~\ref{s:inverse-budget-specialization}; the proof here retains
the alternative calculation with an ordered product of inverse densities.

\subsection{Sparse path forests}


Let \(k=m-\lambda_1\), using \(\lambda_i\) for row lengths. For an absolute \(\epsilon>0\) fixed at \(1/1000\), all sufficiently large \(m\), and \(1\le k\le m^{1-\epsilon}\), we have
\begin{equation}
 J_m(\lambda,p)\le m^{-10k}
 \label{ten:inverse:eq:1}
\end{equation}
for \(p\) larger than an absolute constant. Here are details of a contraction giving this estimate.

Use ordered \(k\)-tuples of distinct positions, containing the representation \(\lambda\) by branching (and in particular \(D_\lambda\le m^k\)). Here this module is induced from the trivial representation of \(S_{m-k}\), since that is the stabilizer. In this module every copy of \(\lambda\) is in the orthocomplement of the subspaces of functions omitting any one coordinate: the \((k-1)\)-tuple module has no copy. Thus such functions sum to zero over each coordinate given all others. In a one-sweep transition between distinct tuples \(x,y\), use the candidate paths on the cube that take their successive updated bits from \(y\), starting at \(x\). For each pair \(i,j\), let \(h_{ij}=0\) unless their paths at an update are at opposite sites of a common gate (matched pair), at the axis of their last starting-bit difference in update order. In the latter case set \(h_{ij}=1\) if they take opposite exits and \(-1\) if they take the same exit. The transition probability is
\[
 m^{-k}F_{[k]},\qquad F_L=\prod_{i<j\ {\rm in}\ L}(1+h_{ij}).
\]
Indeed any path conflict must first arise at a common gate with different entry sites, then necessarily at the indicated axis for that pair if this is their first conflict. With no conflicts, each gate used twice saves a factor of two in probability. Also \(m^{-k}F_L\) is the path probability for true sweep paths on coordinates of \(L\) with independent fair-bit paths on all others (restricted here to distinct endpoints).

For the action on zero-sum functions as above we may replace \(F_{[k]}\) in the expectation kernel (from earlier \(x\) to later \(y\)) by \(G=\sum_{L\subseteq[k]}(-1)^{k-|L|}F_L\), since for every proper subset the sum against such a function in \(y\) vanishes by summing a missing coordinate first. This alternating sum cancels by toggling inclusion of any isolated vertex, so it is zero unless in the interaction graph (\(h_{ij}\ne0\) edges) there are no isolated vertices. Use this replacement at the later of two successive sweeps when taking the expectation of a function of the final tuple. The orders of the two sweeps need not agree. For the sup norm contraction it suffices to bound, uniformly in the start and \(L\), the probability of this covering event after the first true sweep and then the \(F_L\) experiment for the second sweep, paying a factor \(2^k\). We need only upper bound the event, so can drop the distinct-endpoint requirement of that experiment, testing if every path shared some gate with another.

After the first sweep the random occupied \(k\)-set \(U\) satisfies
\[
 \Pr(V_0\subseteq U)\le (k/m)^{|V_0|}
\]
for each set \(V_0\). In fact the inequality bounding joint inclusion probability by the product of one-site marginals holds starting from a fixed set and persists after independent fair matching switches: average the old product bound over the switches on pairs. On a fully tested pair the product of old marginals is at most the square of their mean. And the sweep makes the one-site marginals constant.

Generate the following graph independently of \(U\) on all \(m\) sites. At each site start two potential paths, one using a common random sweep switching system, the other having its own independent random fair-bit path. Join sites if any two paths (one from each) share a gate at the same update, allowing any choices of their two types. Regardless of the coordinate labeling of \(U\), the second experiment can use a path of the requested type at each site, so our event requires \(U\) to induce a graph with no isolates.

Such a graph contains a spanning forest without singleton components. For a fixed forest on \(k\) indices with \(j\) edges, the expected number of \emph{ordered distinct} site placements realizing all its edges is at most \(m^k(8\log_2 m/m)^j\).

To count, condition on the common switching system, sum over which of the two paths at each endpoint and which layer on each edge witnesses the encounter, and upper bound distinct placements by iid uniform placements times \(m^k\), using separate independent tapes for the private type per index (correct for distinct placements). Before imposing edge constraints the individual two-type records of positions are then independent across indices and each path has uniform one-site marginals. Peeling leaves gives probability at most \((2/m)^j\) for those constraints. Using the joint inclusion bound and dividing the labeled count (summed over forests) by \(k!\), since a successful induced graph has some witnessing forest for \emph{each} indexing of the occupied sites, the required sup norm contraction is bounded by
\[
 2^k \frac{k^k}{k!}\sum_{\lceil k/2\rceil\le j\le k-1}
 \binom{\binom{k}{2}}{j}(8\log_2 m/m)^j.
\]
For the empty sum use zero. Use one forward and one adjoint sweep so this bounds the spectral radius of \(B_m^*B_m\) on the indicated isotypic space by the sup norm estimate (permutation action on functions, or expectations for the self-adjoint kernel). For large \(m\) in the range of \eqref{ten:inverse:eq:1} the bound is at most \(m^{-\epsilon k/3}\) since \(\binom{\binom{k}{2}}j\le(O(k))^j\). Thus the singular values there are at most \(m^{-\epsilon k/6}\), proving \eqref{ten:inverse:eq:1} by the dimension bound.

Two other dimension facts we use are as follows. If \(\lambda\) has more than \(m/2\) rows, even one axis update annihilates it: matching-swaps invariants occur only for diagrams dominating the partition made from \(m/2\) parts of size 2, by Young's rule. For remaining \(\lambda\), if \(k>m^{1-\epsilon}\), then \(H_\lambda\gtrsim m^{1-\epsilon}\) for large \(m\). Indeed if \(k\le m/4\), the hook formula gives \(D_\lambda\ge e^{-1}\binom mk\): the leg corrections to first-row hook lengths inflate their product over \(\lambda_1!\) by at most \(\exp(k/(\lambda_1-\lambda_2+1))\). Otherwise width and height are both at most \(3m/4\). If their sum is at most \(m/(2e)\), the hook formula suffices. Else transpose if necessary to have first row \(s\ge m/(4e)\), and restrict to a subshape with first row \(s\) and another \(\lfloor s/4\rfloor\) cells, using the previous first-row argument and branching.

\subsection{The spin and marked-site estimate}


\begin{proposition}[Marked-grid moment estimate]
\label{ten:inverse:main}
Say that a diagram fits the \(q\)-hook if all cells belong to its first \(q\) rows or its first \(q\) columns. Write \(e_m(l)=l\log(m/l)\), with \(e_m(0)=0\). Fix \(\theta=1/2\), \(c_0=1/100\). Consider \(n=ab\ge2\), where \(a,b\) are powers of two with \(\log_2 a,\log_2 b\) differing by at most one, and \(q\le n^{1/16}\) a positive integer. Let \(p\ge2\) be real and let \(\mathcal B_a,\mathcal B_b\ge0\). Suppose that, for each \(m\in\{a,b\}\), every diagram \(\xi\vdash m\), and every \(q\)-hook subdiagram \(\eta\subseteq\xi\) of size \(m-l\),
\begin{equation}
 \log J_m(\xi,p)\le \mathcal B_m+\theta H_\eta+c_0\log(m)l-e_m(l),
 \label{ten:inverse:eq:2}
\end{equation}
using \(-\infty\) for log zero. Then for every \(\lambda\vdash n\) and every \(q\)-hook subdiagram \(\omega\subseteq\lambda\) of size \(g=n-t\),
\begin{equation}
 \log J_n(\lambda,p)
 \le \theta H_\omega+c_0\log(n)t-e_n(t)+
 a\mathcal B_b+b\mathcal B_a+O\big(t+(a+b)(q^2+n^{1/4})\log(n+1)\big),
 \label{ten:inverse:eq:3}
\end{equation}
The implicit constant is absolute.

\end{proposition}

If one grid factor has size \(1\), balance and \(n\ge2\) force \(n=2\),
and the other factor has size \(m=n\). Apply \eqref{ten:inverse:eq:2}
to that factor with \((\xi,\eta,l)=(\lambda,\omega,t)\). The allowance
in \eqref{ten:inverse:eq:3} is \(\mathcal B_2+2\mathcal B_1\), so the
conclusion already holds with zero error because \(\mathcal B_1\ge0\).
We may now assume \(a,b\ge2\).

To prove this, in the irreducible \(\lambda\) bound \({\rm Tr}_\lambda|CR|^p\) by \({\rm Tr}_\lambda XY\), where \(X=|R^*|^p,\ Y=|C|^p\), by Araki-Lieb-Thirring and cyclic invariance for the trace of powers. These are positive subgroup elements, each a tensor product. They may each first be split into a sum by the diagrams \(\xi\) on the local symmetric groups (compressions by central subgroup isotypic projections). Use \(X,Y\) also for a fixed resulting pair of terms, still positive group algebra elements.

Use an induced module of cell states as follows. There are \(t\) distinct labeled holes, one site each, and the other \(g\) sites carry spins with value space \(V=V_+\oplus V_-\) (dimensions \(q,q\)). Site permutations act by permuting and using the sign on the odd spins, i.e. a relative-order sign on the \(V_-\) factors in the usual graded tensor permutation. Holes are even, without signs. Equivalently one takes a direct sum over hole positions and plus/minus patterns, transporting the spin tensors by permutations with that sign. Projection to one placement \(z\) of labeled holes is denoted \(D_z\), leaving free spins. In products computed with reordered site factors the graded reordering is a unitary identification; all subgroup operations within disjoint blocks preserve grading and act as usual tensor products after grouping their sites.

With chosen plus and minus counts and local Schur-Weyl types \(\alpha,\beta\) for the two spin spaces, spectra on these types correspond to highest weights \(\alpha,\beta\), both of length \(\le q\). For \(s\) spins, the permutation module on this sector consists (apart from the general-linear representation spaces) of induction from the symmetric-group types \(\alpha,\beta'\), with prime denoting transpose, by the odd-factor sign. In the presence of \(l\) additional labeled holes the whole representation induces further from \(S_s\) to \(S_{s+l}\).

Choose a global sector of this kind with plus type \(\mu\) made of the first \(q\) rows of \(\omega\) and minus type \(\nu\) the transpose of the remaining rows. Write \(Z\) for its orthogonal projector (full isotypic under the block-unitary or equivalently general-linear action). The spin module within \(g\) sites contains \(\omega\), since \(c_{\mu,\nu'}^\omega\ge1\) by the LR rule (the skew part \(\omega/\mu\) just consists of the lower rows). With the holes the sector thus has \(\lambda\)-multiplicity at least \(f^{\lambda/\omega}\) by branching. Also
\[
 D_\lambda\le \binom nt D_\omega f^{\lambda/\omega}
\]
by splitting standard tableau fillings into these two cell sets. Let \(P_\lambda\) denote the global group-type projector.

We detail the trace estimate on this module. For any block-diagonal positive definite spin density \(D\) of trace 1 (blocks for \(+,-\)), put \(h=D^{-1}\), with \(h_{[g]}\) acting by \(h\) on all spins, same matrix independent of which sites, and identity transport (no change) on hole locations. On the sector \(Z\), \(h_{[g]}\) has minimum eigenvalue \(\ge\exp(J)\), where
\[
 J=g\,\mathsf H\big((\mu_i/g)_i,(\nu_j/g)_j\big)
\]
and \(\mathsf H\) is Shannon entropy in natural logarithms (use \(J=0\) for \(g=0\)). In fact weights in each general-linear representation are majorized by its highest weight, so the maximum eigenvalue of the density power on the sector is bounded by its product of sorted block eigenvalues raised to highest-weight powers, which is at most \(\exp(-J)\) by the entropy/product inequality with eigenvalues summing in total to 1. The multinomial bound gives \(e^J\ge\binom g{|\mu|} D_\mu D_\nu\ge D_\omega\). Thus with \(u=1-\theta\),
\begin{equation}
 ZP_\lambda\le D_\omega^{-u} h_{[g]}^u.
 \label{ten:inverse:eq:4}
\end{equation}
Both sides commute with group algebra elements.

\subsection{Local polynomial dimensions}
Some local polynomial factors used below are recalled here. On \(s\le m\) spins of plus/minus types \(\alpha,\beta\), an occurring permutation type \(\eta\) has
\begin{equation}
 D_\eta \ \ge\ (m+q+1)^{-O(q^2)}
 \exp\big(s\,\mathsf H(\alpha/s,\beta/s)\big).
 \label{ten:inverse:eq:5}
\end{equation}
Indeed LR implies \(\eta_i\le\alpha_i+q,\ \eta'_j\le\beta_j+q\), and \(\eta\) fits the \(q\)-hook. For the first inequality a row of an LR tableau with content \(\beta'\) has at most \(\beta'_1\) boxes: match each entry \(>1\) to a distinct preceding entry of value one less by the lattice word property. Predecessors must be in strictly earlier rows, so the chains ending for boxes of this row use distinct ones. Transposition gives the other inequality. The product of hook lengths off the \(q\times q\) square in the first \(q\) rows is at most \((s+q+1)^{O(q^2)}\prod\alpha_i!\): the leg lengths there are at most \(q\), and row portions there have lengths at most \(\alpha_i\). Use the analogous column estimate, and hooks \(\le s\) on the square. The hook formula and multinomial estimate now give \eqref{ten:inverse:eq:5}, trivially also for \(s=0\). Consequently the multiplicity of \(\eta\) in this sector is at most \((m+q+1)^{O(q^2)}\) since its LR coefficient is at most \(\binom{s}{|\alpha|}D_\alpha D_\beta / D_\eta\), and general-linear dimensions satisfy this polynomial-factor bound by the Weyl formula.

\subsection{Twirl the rows before splitting the columns}
Split \(X\) additionally by projecting orthogonally according to the number \(l_i\) of holes per row \(i\), and by plus/minus counts and general-linear types of spins there, and symmetric-group type \(\eta_i\) on the nonhole sites there. For the latter use a direct sum of the spin permutation-type projectors over hole placements. All local sectors are equivariant under the rows subgroup, and \(X\) preserves them, giving positive terms. Whenever nonzero, \(\eta_i\) is contained in the previously specified total row type \(\xi_i\) by branching: within each assignment of labels to blocks, a block's sector over all placements is induced from its sector of stabilizer representation at one placement. Before using \eqref{ten:inverse:eq:4} we perform a twirl decomposition of each positive row-sector term. For row \(i\) with \(s_i=b-l_i>0\) and spin types \(\alpha_i,\beta_i\), take a random block-unitarily conjugated density \(\rho_i\) on \(V\) with block spectra \(\alpha_i/s_i,\beta_i/s_i\) (zeros allowed), independent Haar rotations. The average of its spin tensor power with exponent \(u\) per density (using \((\rho_i^u)^{\otimes s_i}\)) acts scalarly on the specified sector with scalar at least
\[
 (n+1)^{-O(q^2)}\exp(-u\, s_i\mathsf H(\alpha_i/s_i,\beta_i/s_i)).
\]
This follows by Haar averaging, taking at least the highest-weight eigenvalue divided by the general-linear dimensions. Indeed the power acts by the general-linear representation, identically in every multiplicity copy, and its conjugation average over \(U(q)\times U(q)\) replaces each general-linear irrep matrix by a scalar by Schur's lemma. For an empty spin set take scalar 1 and any density. The row term itself commutes with all such powers in its own row. Thus it is an average, with normalization the reciprocal scalar, of positive terms obtained by sandwiching it with the spin powers having exponent \(u/2\). We use the product over rows, and call the normalized integrand \(X'\), so the normalization consists of the product of those reciprocal scalars. Use positive definite approximations, if needed, at factors tending to 1 in the estimates by adding small positive scalar matrices to the densities and renormalizing, then taking limits.

Choose $D=\sum_i\rho_i/a$ in \eqref{ten:inverse:eq:4}. \paragraph{The comparison before the column split.}
Put $A=ZP_\lambda$, $H=h_{[g]}^u$, and $c=D_\omega^{-u}$, so that
$A\le cH$.  At this stage $Y$ is still a positive element of the column
group algebra, possibly split by central total-column types.  Thus
$Y$ commutes with both $A$ and $H$.  For every positive row-twirl term
$X'$, these commutations give
\[
 c\operatorname{Tr}(X'YH)-\operatorname{Tr}(X'YA)
 =\operatorname{Tr}\!\left(
 X'Y^{1/2}(cH-A)Y^{1/2}\right)\ge0.
\]
The last trace is nonnegative because its two factors are positive.
Moreover
\[
 \operatorname{Tr}(X'YH)=\operatorname{Tr}(H^{1/2}X'H^{1/2}Y).
\]  No commutation of $X'$ with
$A$ or $H$ has been used.  We now split $Y$ into its positive local
spin and hole sectors and apply the column twirl.  Performing that split
before the comparison would require commutations that its summands need
not satisfy.


Only now split \(Y\) by the analogous hole and spin sectors on columns and twirl with column densities \(\sigma_j\) at exponent \(u\) as above. The right-hand trace needs only the diagonal compressions of the two terms (row and column) to \(D_z\), summed over placements: a row operator preserves the row of each hole label, and a column operator its column, so there is only the original placement as intermediate placement in this trace.

\subsection{Positive compression with labeled holes}
On a single local block of size \(m\), with \(l\) holes at specified sites, write \(\xi,\eta,\alpha,\beta\) for its total permutation type, nonhole spin permutation type, and spin general-linear types. Before inserting \(h\) and before the spin-power sandwiching and normalization, the compressed PSD operator in the selected sector has trace at most
\begin{equation}
 (m+q+1)^{O(q^2)}
 \frac{J_m(\xi,p)}{(m)_l D_\eta}.
 \label{ten:inverse:eq:6}
\end{equation}
Here \((m)_l=m!/(m-l)!\). To see this, taking the diagonal with labeled holes fixed retains only coefficients in the stabilizer \(S_{m-l}\) of those positions from the local group algebra term of \(X\) or \(Y\) on type \(\xi\). Here we can identify the local module independently for each assigned set of hole labels, grouping the sites blockwise using the graded permutation identification. On fixed placements the action of the retained coefficients is thus via the spin permutation action in this block. The coefficient restriction remains positive (conditional expectation in group algebra, or compression of the regular matrix). Its identity coefficient is \(J_m(\xi,p)/m!\), unchanged, so by regular trace on \(S_{m-l}\) its trace on type \(\eta\) is at most \(J_m(\xi,p)/((m)_l D_\eta)\). Now apply the local multiplicity bound. All sector compressions here besides the original total-type restriction use projectors preserving fixed placements and commuting there with the retained subgroup action.

\subsection{The interpolated density product}
For clarity, discard terms with a zero bound and divide out the product of the trace bounds \eqref{ten:inverse:eq:6} over blocks on each respective side, including polynomial factors, so the remaining spin operators on each side prior to the sandwiches are positive \(X_0,Y_0\) with trace and norm at most 1 (formed locally for the spin system for \(z\), with graded grouping). In terms of tensor densities on occupied spin sites formed from the row matrices and the column matrices, denoted by \(\rho_R,\sigma_C\), and \(h_* = h^{\otimes g}\), the needed extra trace factor is bounded by
\begin{equation}
 \left\| X_0^{1/2}\rho_R^{u/2} h_*^{u/2}\sigma_C^{u/2}Y_0^{1/2}\right\|_{\rm HS}^2
 \le e^t .
 \label{ten:inverse:eq:7}
\end{equation}
\paragraph{Interpolation in the original matrix order.}
Fix the occupied cells $\mathcal E\subseteq[a]\times[b]$, with
$|\mathcal E|=g=ab-t$, and use one fixed ordering of them.  For the row
and column trace-one spin densities $\rho_i,\sigma_j$, write
\[
 D=\frac1a\sum_i\rho_i,\quad h=D^{-1},\qquad
 \rho_R=\bigotimes_{(i,j)\in\mathcal E}\rho_i,\quad
 \sigma_C=\bigotimes_{(i,j)\in\mathcal E}\sigma_j,\quad
 h_*=h^{\otimes g}.
\]
For the moment assume the densities are positive definite.  The normalized
unsandwiched compressions $X_0,Y_0$ are positive with traces at most one;
in particular they are contractions.  On the strip $0\le\Re z\le1$,
consider the Hilbert--Schmidt-valued analytic function
\[
 F(z)=X_0^{1/2}\rho_R^{z/2}h_*^{z/2}\sigma_C^{z/2}Y_0^{1/2}.
\]
On the lower boundary, the product of the three imaginary powers is
unitary, and hence
\[
 \|F(iv)\|_{\mathrm{HS}}
 \le\|X_0^{1/2}\|_{\mathrm{HS}}\|Y_0^{1/2}\|_{\mathrm{op}}\le1.
\]
On the upper boundary there is the exact factorization
\[
 F(1+iv)=
 \bigl[X_0^{1/2}\rho_R^{iv/2}\bigr]
 \bigl[\rho_R^{1/2}h_*^{1/2}h_*^{iv/2}\sigma_C^{1/2}\bigr]
 \bigl[\sigma_C^{iv/2}Y_0^{1/2}\bigr].
\]
Removing the two outer contractions in Hilbert--Schmidt norm and
factorizing the tensor trace yields
\[
 \|F(1+iv)\|_{\mathrm{HS}}^2
 \le\prod_{(i,j)\in\mathcal E}
 \operatorname{Tr}\!\left(h^{1/2}\rho_i h^{1/2}
                                  \widetilde\sigma_j\right),
 \qquad
 \widetilde\sigma_j=h^{iv/2}\sigma_jh^{-iv/2}.
\]
Every $\widetilde\sigma_j$ is a trace-one positive density.  The same
unitary conjugation occurs in every row, so each full column, including
its omitted cells, has mean
\[
 \frac1a\sum_i
 \operatorname{Tr}\!\left(h^{1/2}\rho_i h^{1/2}
                                  \widetilde\sigma_j\right)
 =\operatorname{Tr}\widetilde\sigma_j=1.
\]
If column $j$ has $l'_j$ holes, arithmetic--geometric means bounds its
occupied product by
\[
 \left(\frac{a}{a-l'_j}\right)^{a-l'_j}\le e^{l'_j},
\]
with the empty product equal to one.  Thus
$\|F(1+iv)\|_{\mathrm{HS}}\le e^{t/2}$.  Applying the three-lines
theorem to the scalar function $\operatorname{Tr}(T^*F(z))$ for every
Hilbert--Schmidt unit vector $T$ gives
\[
 \|F(u)\|_{\mathrm{HS}}^2\le e^{ut}\le e^t,
 \qquad 0\le u\le1.
\]
The trace after the two density sandwiches is exactly this squared norm:
\begin{align*}
 &\operatorname{Tr}\!\left[
 h_*^{u/2}(\rho_R^{u/2}X_0\rho_R^{u/2})h_*^{u/2}
            (\sigma_C^{u/2}Y_0\sigma_C^{u/2})\right]
   =\|F(u)\|_{\mathrm{HS}}^2.
\end{align*}
Only cyclic invariance of trace is needed for this identity.  At no point
have distinct positive matrices been commuted.

For singular densities, replace each $\rho_i$ and $\sigma_j$ by
$(\rho_i+\varepsilon I)/(1+2q\varepsilon)$ and
$(\sigma_j+\varepsilon I)/(1+2q\varepsilon)$.  Keep
$\varepsilon>0$ throughout the twirls and the trace comparison.  The
bounds above are uniform in $\varepsilon$; eliminate the inverse density
using them before letting $\varepsilon$ decrease to zero in the twirl
scalars.  This procedure remains valid when the limiting row mean is
singular and does not require an inverse at that limit.  All the spin
matrices preserve parity, so graded regrouping conjugates the whole
calculation by unitaries and changes neither positivity nor these norms.

This proves \eqref{ten:inverse:eq:7}.

By \eqref{ten:inverse:eq:5}, the reciprocal twirl scalars are bounded by local factors \(D_\eta^u\) times \((n+1)^{O(q^2)}\) (in the regularized case multiply by factors tending to 1). Combined with \eqref{ten:inverse:eq:6} this leaves, apart from polynomial factors and \eqref{ten:inverse:eq:7}, one factor
\[
 J_m(\xi,p) D_\eta^{-\theta}/(m)_l
\]
per block on both sides. The resulting bound for \({\rm Tr}(P_\lambda Z XY)\) carries also the factor \(D_\omega^{-u}\). To get a bound on \(D_\lambda{\rm Tr}_\lambda XY\) pay at most \(\binom nt D_\omega\) by the dimension and multiplicity estimates above. Traces before expansion for \(P_\lambda Z\) give exactly the sector multiplicity times the irrep trace.

\subsection{Entropy accounting for the marks}
For fixed vectors of hole counts with sums \(t\) on each side, the number of placements is at most \((t!/\prod l_i!)(t!/\prod l'_j!)\). Across both sides the product of \(1/(m)_l\) is at most \(n^{-t}e^{2t}\), using \((m)_l\ge (m/e)^l\). We record the entropy accounting: on the row side,
\[
 \log(t!/\prod l_i!)\le t\log a+\sum_i e_b(l_i)-e_n(t)
\]
and analogously on columns. Thus the \(\sum e_m(l)\) terms cancel those of \eqref{ten:inverse:eq:2}, and \(\binom nt\le\exp(e_n(t)+t)\). Also \(\sum l\log m=t\log n\). These estimates apply including zero counts. Finally the number of count and type choices and the polynomial factors combined are bounded by \(\exp(O((a+b)(q^2+n^{1/4})\log(n+1)))\), using at most \((m+1)^{O(\sqrt m)}\) diagrams of size \(\le m\) (e.g. describe by Durfee square and row/column lengths along it). There is no extra sum over hole-label allocations besides placements. This proves \eqref{ten:inverse:eq:3}.

\subsection{Extending the child estimate to an inherited hook}
In this subsection $E_m(\xi;p)=\log J_m(\xi,p)$ and
$j_m(l)=e_m(l)$, as in the marked recurrence.
\label{s:child-hook-extension-section}

The parent uses its own hook parameter when invoking the marked recurrence.
It is therefore necessary to extend the child induction from the child's
smaller hook to every hook inherited from the parent.  The extension below
in fact applies to every subpartition of the child diagram.

Use $\epsilon,c_0,\theta,Q_m$ and $G_m$ from
Proposition~\ref{s:inverse-bounded-exponent}.

\begin{lemma}[Allowance for arbitrary inherited subdiagrams]
\label{s:arbitrary-hook-child-allowance}
There are absolute $C$ and $m_0$ with the following property.  Let
$m\ge m_0$ be a power of two and let $p\ge2$.  Suppose
\[
 E_m(\xi;p)\le G_m(\xi)\quad\hbox{for large-level $\xi$},
 \qquad
 E_m(\xi;p)\le0\quad\hbox{for the other $\xi$}.
\]
Then, for every $\xi\vdash m$ and every subpartition
$\eta\subseteq\xi$, putting $l=m-|\eta|$ gives
\begin{equation}
 E_m(\xi;p)
 \le \mathcal B_m+\theta H_\eta+c_0l\log m-j_m(l),
 \qquad \mathcal B_m=Cm^{1-4\epsilon}\log(m+1).
 \label{s:eq:arbitrary-hook-child-allowance}
\end{equation}
This includes every $q$-hook subpartition, for any inherited $q\ge1$.
The constants do not depend on $p$ or $q$.
\end{lemma}
\begin{proof}
A zero moment causes no difficulty.  If $\xi$ is not large-level, use
$H_\eta\ge0$ and the elementary minimum
\[
 \min_{0\le x\le m}
 \bigl(c_0x\log m-j_m(x)\bigr)
 =-e^{-1}m^{1-c_0}.
\]
Since $c_0>4\epsilon$, this negative part is bounded in absolute value
by the stated $\mathcal B_m$.  This also handles the trivial diagram, whose
moment is exactly one.

Suppose next that $\xi$ is large-level.  Intersect $\eta$ with the
$Q_m$-hook to obtain $\alpha$, and let
$h=|\eta|-|\alpha|$.  Since $\alpha\subseteq\xi$, the hypothesis gives
\[
 E_m(\xi;p)
 \le\theta H_\alpha+c_0(l+h)\log m-j_m(l+h).
\]
The derivative $j_m'(x)=\log(m/x)-1$ is at least $-1$ on $(0,m]$.
By continuity at zero,
\[
 j_m(l)-j_m(l+h)\le h.
\]
Consequently the excess over the claimed expression without $\mathcal B_m$ is
at most
\[
 c_0h\log m+h-\theta(H_\eta-H_\alpha).
 \tag{$\dagger$}
\]
Tableau extension gives $H_\eta\ge H_\alpha$.  If
$h<m^{1-4\epsilon}$, dropping the last term in $(\dagger)$ proves the
required allowance.

For larger $h$, the deleted region is a translated straight partition
$\beta$: its row lengths are
$(\eta_{Q_m+i}-Q_m)_+$, $i\ge1$.  Both its width and its height are at
most $m/Q_m$, so every hook length of $\beta$ is at most $2m/Q_m$.
The hook formula and $h!\ge(h/e)^h$ imply
\[
 \log D_\beta\ge h\log\frac{Q_mh}{2em}.
\]
Fill $\alpha$ with the first $|\alpha|$ entries of a tableau and fill
the translated $\beta$ with the remaining entries.  The boundary
inequalities hold because $\alpha$ is an order ideal in $\eta$, giving
$D_\eta\ge D_\alpha D_\beta$.  Therefore
\[
 H_\eta-H_\alpha
 \ge h\log\frac{Q_mh}{2em}
 \ge h\left[\left(\frac1{16}-4\epsilon\right)\log m-\log(4e)\right]
\]
for sufficiently large $m$, when $Q_m\ge m^{1/16}/2$.  The numerical
inequality
\[
 \theta\left(\frac1{16}-4\epsilon\right)
 =0.02925>0.01=c_0
\]
shows that $(\dagger)$ is nonpositive once $m\ge m_0$.  This proves the
lemma.
\end{proof}

If the two children have different already established exponents, use
their maximum $p$.  Each sweep is a contraction, so increasing the real
Schatten exponent only decreases every $E_m(\xi;p)$.  Apply the lemma
with that common $p$ and then the marked recurrence with the parent
parameter $q=Q_n$.  This verifies the all-subdiagram quantifier in the
child hypothesis; restricting it to $Q_m$-hook diagrams would not suffice.

\subsection{The bounded-exponent induction}


\begin{proof}[Proof of Proposition~\ref{s:inverse-bounded-exponent}]
The trivial diagram has moment one and is not large-level. A zero moment
satisfies either branch. Once the finite starting range has been fixed,
Lemma~\ref{lem:finitegap} permits a common exponent large enough for
every other diagram in that range, including a possibly negative
$G_n(\lambda)$. It will also be chosen above the fixed sparse
requirement in \eqref{ten:inverse:eq:1}. The sparse branch then holds
whenever that estimate applies; its remaining finitely many sizes will
belong to the starting range. We now fix the large-size thresholds
without using that exponent.

For larger \(d\), split the dimensions into \(\lfloor d/2\rfloor,\lceil d/2\rceil\) yielding \(a,b\). Let \(p\) be the maximum of the child exponents. Lemma~\ref{s:arbitrary-hook-child-allowance} gives
\eqref{ten:inverse:eq:2} for every parent-hook subdiagram with
$\mathcal B_m=O(m^{1-4\epsilon}\log(m+1))$. Increasing the child exponent to
their maximum is allowed by contraction. Thus the displayed all-subdiagram
hypothesis, including the inherited parent hook, is satisfied.

We next check three comparisons needed to apply
\eqref{ten:inverse:eq:3}. We may assume the parent is large-level
and has positive moment, so the sparse dimension estimate gives
$H=H_\lambda\ge c n^{1-\epsilon}$.

First,
\[
 G_n(\lambda)\le0.6H
\]
for all sufficiently large $n$. To see this, test the canonical
$Q_n$-hook portion $\alpha\subseteq\lambda$, and put $s=n-|\alpha|$.
If $s\ge n^{1-4\epsilon}$, the hook calculation in the preceding
lemma gives
\[
 H-H_\alpha\ge
 s\left[\left(\frac1{16}-4\epsilon\right)\log n-\log(4e)\right].
\]
Multiplied by $\theta$, this is larger than $c_0s\log n$.
Consequently this candidate costs at most $\theta H$, even after
its negative deletion entropy is discarded. If $s<n^{1-4\epsilon}$,
the candidate costs at most
$\theta H+c_0n^{1-4\epsilon}\log n=\theta H+o(H)$.
Both cases prove the asserted bound.

Second, the deletion count $t$ of an actual minimizing $\omega$
satisfies
\[
 t=O\left(n^{1-4\epsilon}+\frac H{\log n}\right).
\]
Indeed if $t\ge n^{1-4\epsilon}$, then
$\log(n/t)\le4\epsilon\log n$ and $H_\omega\ge0$, so
\[
 (c_0-4\epsilon)t\log n\le G_n(\lambda)\le0.6H.
\]
Below that threshold the displayed size bound is immediate.

Third, use this minimizing $\omega$ in \eqref{ten:inverse:eq:3}.
The child allowances and the grid error satisfy
\[
 \begin{aligned}
 a\mathcal B_b+b\mathcal B_a
   &=O(n^{1-2\epsilon}\log(n+1)),\\
 (a+b)(Q_n^2+n^{1/4})\log(n+1)
   &=O(n^{3/4}\log(n+1)).
 \end{aligned}
\]
Together with the bound on $t$ and $H\ge c n^{1-\epsilon}$,
these imply, for an absolute $K$ independent of $p$,
\[
 \log J_n(\lambda,p)\le G_n(\lambda)+KH/d.
\]
For example the first error divided by $H/d$ is
$O(d n^{-\epsilon}\log(n+1))=o(1)$; the other sublinear powers
are smaller. The $H/\log n$ part of the deletion count is already
$O(H/d)$.

For a positive moment and $\delta>0$, the elementary inequality
$\sum_i x_i^{1+\delta}\le(\sum_i x_i)^{1+\delta}$ for $x_i\ge0$
gives
\[
 \log J_n(\lambda,p(1+\delta))
 \le(1+\delta)\log J_n(\lambda,p)-\delta H.
\]
Choose the size threshold so that $K/d\le0.1$. The preliminary
bound is then at most $0.7H$, and taking $\delta=A/d$ with
$A\ge4K$ yields
\[
 \log J_n(\lambda,p(1+A/d))
 \le G_n(\lambda)+(K-0.3A)H/d
 \le G_n(\lambda).
\]
All thresholds and $A$ are independent of the starting exponent.
Choose an integer $d_0$ above them, also large enough that every child
in a recursion step is in the range of the inherited-hook lemma.
Now choose one starting exponent $P_*$ for $d\le d_0$ as described
at the beginning of the proof, and set
\[
 p_d=P_*\quad(d\le d_0),\qquad
 p_d=(1+A/d)\max\{p_{\lfloor d/2\rfloor},p_{\lceil d/2\rceil}\}
       \quad(d>d_0).
\]
Strong induction proves both branches of the proposition. To bound
these exponents, follow a child attaining each maximum until the first
index at most $d_0$. Along this path, parent $u$ and child $v$ obey
$u-1\ge2(v-1)$, so the sum of reciprocal bit lengths above the
starting range is at most $2/d_0$. Thus
$p_d\le P_*e^{2A/d_0}$ for every $d$.
\end{proof}

\subsection{From moments to forward sweeps}
Put $P=\sup_d p_d<\infty$ and choose an integer $r\ge4P$. At large levels, the bound $G_n(\lambda)\le0.6H_\lambda$ and the moment-slack inequality give
\[
 \log J_n(\lambda,2r)
 \le\left(1-0.4\frac{2r}{p_d}\right)H_\lambda
 \le-2H_\lambda.
\]
Consequently
\[
 \sum_{\lambda\ne(n)} J_n(\lambda,2r)=o(1).
\]
The large-level bound is summable with negligible total by the partition count and dimension bound. At the others use \eqref{ten:inverse:eq:1} (at most \(2^k\) diagrams per sparse \(k\)). By Schatten H\"older, the analogous regular-representation Hilbert-Schmidt sum with \({\rm Tr}_\lambda|B_n^r|^2\) in place of \({\rm Tr}_\lambda|B_n|^{2r}\) is no larger. This sum excluding the trivial diagram is the squared \(L^2\) distance of the permutation density from 1 relative to uniform, by the regular action trace formula. By Cauchy-Schwarz, \(rd\) shuffles therefore suffice in total variation at large \(d\), uniformly over starting orders by relabeling.
